% 2-5 Lesson Quiz, short form (from the 2-5 Lesson Quiz block in a2_2-5_TE.tex).
% Kept: Savvas items 2, 3, 5 (numbered as on the Savvas quiz so they match the packet's "On the quiz" box).
% Left off: items 1 and 4. No item on this quiz has a graph.
\documentclass[11pt]{article}
\usepackage[margin=0.7in]{geometry}
\usepackage{amsmath,amssymb}
\usepackage{enumitem}
\usepackage{tikz}
\usepackage{pgfplots}
\pgfplotsset{compat=1.18}
\usepackage{fancyhdr}
\usepackage{parskip}
\setlist[enumerate]{leftmargin=*, itemsep=6pt}
\pagestyle{fancy}
\fancyhf{}
\fancyfoot[L]{\small\textit{en}Vision\textsuperscript{\textregistered} Algebra 2}
\fancyfoot[C]{\small 2-5 Lesson Quiz, short form}
\fancyfoot[R]{\small\thepage}
\renewcommand{\headrulewidth}{0pt}
\newcommand{\choice}[1]{\tikz[baseline=-0.6ex]\node[draw, circle, inner sep=0.6pt, minimum size=3.3mm, font=\scriptsize\bfseries] {#1};\;}
\newcommand{\checkbox}{$\square$\;}
\newcommand{\ansbox}{\fbox{\rule{0pt}{16pt}\hspace{1.2cm}}}
\pgfplotsset{quizgrid/.style={
  axis lines=middle, grid=both, grid style={gray!45},
  ticklabel style={font=\scriptsize, fill=white, inner sep=1pt},
  xlabel={$x$}, ylabel={$y$}, axis line style={thick}, clip=true}}

\begin{document}

\noindent\textbf{Name}\; \rule{6cm}{0.4pt} \hfill Period \rule{1.2cm}{0.4pt} \quad Date \rule{2.2cm}{0.4pt}

\vspace{4pt}\noindent\rule{\textwidth}{0.6pt}

\vspace{2pt}\noindent{\LARGE\textbf{2-5 \; Lesson Quiz: Completing the Square}} \hfill {\small short form}

\vspace{4pt}\noindent\rule{\textwidth}{0.6pt}

\vspace{2pt}
{\small Question numbers match the Savvas Lesson Quiz. Show your work: the line where you complete the square counts.}

\vspace{6pt}
\begin{minipage}[t]{0.50\textwidth}
\textbf{2.} Write the equation $x^2+24x+60=0$ in the form $(x+p)^2=q$.

\vspace{4pt}
\begin{itemize}[leftmargin=1.2em, itemsep=5pt]
\item[] \choice{A} $(x+12)^2=-60$
\item[] \choice{B} $(x+12)^2=84$
\item[] \choice{C} $(x+12)^2=204$
\item[] \choice{D} $(x+6)^2=-24$
\end{itemize}
\end{minipage}\hfill
\begin{minipage}[t]{0.46\textwidth}
\textbf{Work}

\vspace{4pt}\rule{\linewidth}{0.4pt}

\vspace{14pt}\rule{\linewidth}{0.4pt}

\vspace{14pt}\rule{\linewidth}{0.4pt}
\end{minipage}

\vspace{10pt}\noindent\rule{\textwidth}{0.4pt}

\vspace{6pt}
\begin{minipage}[t]{0.50\textwidth}
\textbf{3.} Solve $0=x^2+6x+13$ by completing the square.

\vspace{4pt}
\begin{itemize}[leftmargin=1.2em, itemsep=5pt]
\item[] \choice{A} $x=5$ and $x=1$
\item[] \choice{B} $x=3+i\sqrt{22}$ and $x=3-i\sqrt{22}$
\item[] \choice{C} $x=-1$ and $x=1$
\item[] \choice{D} $x=-3+2i$ and $x=-3-2i$
\end{itemize}
\end{minipage}\hfill
\begin{minipage}[t]{0.46\textwidth}
\textbf{Work}

\vspace{4pt}\rule{\linewidth}{0.4pt}

\vspace{14pt}\rule{\linewidth}{0.4pt}

\vspace{14pt}\rule{\linewidth}{0.4pt}

\vspace{14pt}\rule{\linewidth}{0.4pt}

\vspace{14pt}\rule{\linewidth}{0.4pt}
\end{minipage}

\vspace{10pt}\noindent\rule{\textwidth}{0.4pt}

\vspace{6pt}
\textbf{5.} Find the zeros of the function $f(x)=-5x^2+4x-1$ by completing the square.

\vspace{6pt}
\begin{minipage}[t]{0.50\textwidth}
\textbf{Work}

\vspace{4pt}\rule{\linewidth}{0.4pt}

\vspace{14pt}\rule{\linewidth}{0.4pt}

\vspace{14pt}\rule{\linewidth}{0.4pt}

\vspace{14pt}\rule{\linewidth}{0.4pt}

\vspace{14pt}\rule{\linewidth}{0.4pt}

\vspace{14pt}\rule{\linewidth}{0.4pt}
\end{minipage}\hfill
\begin{minipage}[t]{0.46\textwidth}
\vspace{18pt}
\centering
$x=\ansbox+\ansbox\,i$ \quad and

\vspace{14pt}
$x=\ansbox-\ansbox\,i$
\end{minipage}

\vfill
\noindent\textbf{Scoring:} Question 3 and Question 5 are 10 points each. Question 2 is a bonus of up to 5 points. Total out of 20.

\newpage
\noindent{\LARGE\textbf{2-5 Lesson Quiz, short form: answer key}}

\vspace{4pt}\noindent\rule{\textwidth}{0.6pt}

\begin{enumerate}
\item[\textbf{2.}] \textbf{B, $(x+12)^2=84$.} $x^2+24x=-60$; $\left(\frac{24}{2}\right)^2=144$; $x^2+24x+144=-60+144$, so $(x+12)^2=84$. Distractors: A adds 144 to the left side only; C adds 144 to $+60$ without moving the constant; D squares half of 12 instead of half of 24. Savvas quiz item 2, DOK 2, skill A30 (Example 2B; the first three lines of Example 3). \textbf{Bonus 5:} 3 for B, 2 for the line that adds 144 to both sides.
\item[\textbf{3.}] \textbf{D, $x=-3+2i$ and $x=-3-2i$.} $x^2+6x=-13$; add $\left(\frac{6}{2}\right)^2=9$ to both sides: $(x+3)^2=-4$; $x+3=\pm2i$; $x=-3\pm2i$. Savvas quiz item 3, DOK 2, skill A30 (Example 3, Try It 3). \textbf{10 points:} 4 for D, 3 for $(x+3)^2=-4$ shown, 3 for $x+3=\pm\sqrt{-4}=\pm2i$ and the final pair.
\item[\textbf{5.}] \textbf{$x=\frac25+\frac15i$ and $x=\frac25-\frac15i$.} Set $f(x)=0$ and divide by $-5$: $x^2-\frac45x+\frac15=0$; $x^2-\frac45x=-\frac15$; add $\left(-\frac25\right)^2=\frac{4}{25}$ to both sides: $\left(x-\frac25\right)^2=-\frac{1}{25}$; $x-\frac25=\pm\frac15i$. Savvas quiz item 5, DOK 2, skill A30 (Example 3 after dividing by $a$; packet Practice block). \textbf{10 points:} 2 for dividing by $-5$, 2 for adding $\frac{4}{25}$ to both sides, 2 for $\left(x-\frac25\right)^2=-\frac1{25}$, 4 for the four boxes (1 each). Accept $0.4\pm0.2i$.
\end{enumerate}

\textbf{Items left off this form} (available for a retake or Form B): Savvas item 1, solve $x^2-18x+81=4$, select all (answer C and D, $x=7$ and $x=11$; Example 1, DOK 2, A30); Savvas item 4, the baking sheet with perimeter 58 in.\ and area 201.25 in.$^2$ (answer A, length 17.5 in.\ and width 11.5 in.; Example 4, DOK 2, A30: $\ell(29-\ell)=201.25$, $(\ell-14.5)^2=9$).

\textbf{IXL replacement:} SmartScore 70 or above on 9MW and NPH replaces this quiz with 20/20.

\end{document}
